\documentclass[11pt,a4paper]{article}

% Template visual Matriz Aprova. Compile com XeLaTeX ou LuaLaTeX.
\usepackage{fontspec}
\usepackage{polyglossia}
\setmainlanguage{portuguese}
\IfFontExistsTF{Aptos}{
  \setmainfont{Aptos}
  \setsansfont{Aptos}
}{
  \IfFontExistsTF{Calibri}{
    \setmainfont{Calibri}
    \setsansfont{Calibri}
  }{
    \IfFontExistsTF{TeX Gyre Heros}{
      \setmainfont{TeX Gyre Heros}
      \setsansfont{TeX Gyre Heros}
    }{
      \setmainfont{Latin Modern Sans}
      \setsansfont{Latin Modern Sans}
    }
  }
}
\usepackage[a4paper,margin=2cm,headheight=16pt]{geometry}
\usepackage{graphicx}
\usepackage{xcolor}
\usepackage{tikz}
\usepackage{tcolorbox}
\usepackage{enumitem}
\usepackage{booktabs}
\usepackage{tabularx}
\usepackage{amssymb}
\usepackage{fancyhdr}
\usepackage{titlesec}
\usepackage{hyperref}

\definecolor{matrizlime}{HTML}{CBFF4D}
\definecolor{matrizink}{HTML}{101313}
\definecolor{matrizmuted}{HTML}{596363}
\definecolor{matrizline}{HTML}{DDE3DF}
\definecolor{matrizsoft}{HTML}{F3F7EC}

\hypersetup{colorlinks=true,linkcolor=matrizink,urlcolor=matrizink}
\pagestyle{fancy}
\fancyhf{}
\lhead{\textsf{\small MATRIZ APROVA}}
\rhead{\textsf{\small APOSTILA}}
\cfoot{\textsf{\small \thepage}}
\renewcommand{\headrulewidth}{0.4pt}
\renewcommand{\headrule}{\hbox to\headwidth{\color{matrizline}\leaders\hrule height \headrulewidth\hfill}}

\titleformat{\section}{\Large\bfseries\color{matrizink}}{\thesection}{0.6em}{}
\titleformat{\subsection}{\large\bfseries\color{matrizink}}{\thesubsection}{0.6em}{}
\setlist{nosep,leftmargin=*}

\newtcolorbox{foco}{colback=matrizsoft,colframe=matrizlime,boxrule=1pt,arc=3pt,left=10pt,right=10pt,top=8pt,bottom=8pt}
\newtcolorbox{lei}{colback=matrizink,colframe=matrizink,coltext=white,boxrule=0pt,arc=3pt,left=10pt,right=10pt,top=8pt,bottom=8pt}
\newtcolorbox{alerta}{colback=white,colframe=matrizline,boxrule=0.6pt,arc=3pt,left=10pt,right=10pt,top=8pt,bottom=8pt}

% Logo usada na capa. Caminho relativo à pasta onde o .tex é compilado
% (materiais/<disciplina>/). Ajuste se a estrutura mudar.
\newcommand{\logomatriz}{../../public/matrizaprova_temaclaro.png}

\begin{document}

\begin{titlepage}
  \thispagestyle{empty}
  \begin{center}
    \includegraphics[width=0.55\linewidth]{\logomatriz}\par
  \end{center}
  \vspace{2.2cm}
  {\sffamily\fontsize{28}{32}\selectfont\bfseries Apostila — Direito Constitucional: Princípios fundamentais\par}
  \vspace{0.4cm}
  {\sffamily\large Receita Federal, CGU, PRF, Polícia Federal, INSS, TRT, TCU/TCE, Câmara e OAB\par}
  \vfill
  \begin{foco}
    \textsf{\textbf{Foco da apostila}}\\[3pt]
    
  \end{foco}
  \vspace{0.7cm}
  {\sffamily\small Atualizado em 16 de agosto de 2026\quad ·\quad Cebraspe, FGV, FCC e similares\par}
  \vspace{1cm}
  {\sffamily\small matrizaprova.com\par}
\end{titlepage}

\tableofcontents
\newpage

% Conteúdo gerado entra aqui.
\section*{O que cai}

Os princípios fundamentais aparecem com frequência em questões literais, associações entre conceitos e situações envolvendo soberania, cidadania, separação de Poderes, objetivos da República e relações internacionais.

Pontos de maior incidência:

\begin{enumerate}
  \item fundamentos da República Federativa do Brasil;
  \item diferença entre fundamentos, objetivos e princípios internacionais;
  \item titularidade e exercício do poder;
  \item separação e independência dos Poderes;
  \item objetivos fundamentais do art. 3º;
  \item princípios que regem o Brasil nas relações internacionais.
\end{enumerate}

\textbf{Atenção:} a banca costuma trocar a categoria do princípio. Cidadania é fundamento; erradicar a pobreza é objetivo; prevalência dos direitos humanos é princípio internacional.

\section*{Teoria essencial}

\subsection*{1. Estado brasileiro e fundamentos da República}

O art. 1º afirma que a República Federativa do Brasil, formada pela união indissolúvel dos Estados, Municípios e Distrito Federal, constitui-se em Estado Democrático de Direito.

Essa frase reúne características importantes:

\begin{itemize}
  \item \textbf{República:} forma de governo baseada, entre outros elementos, na temporariedade e na responsabilidade dos governantes.
  \item \textbf{Federação:} forma de Estado em que há repartição constitucional de competências entre entes autônomos.
  \item \textbf{Estado Democrático de Direito:} o poder é limitado pela Constituição e a legitimidade estatal depende da participação popular e da proteção de direitos fundamentais.
  \item \textbf{União indissolúvel:} não existe direito de secessão dos entes federados.
\end{itemize}

O parágrafo único do art. 1º estabelece que \textbf{todo o poder emana do povo}, que o exerce por meio de representantes eleitos ou diretamente, nos termos da Constituição. O Brasil adota, portanto, democracia representativa com instrumentos de participação direta.

\subsection*{2. Os cinco fundamentos: SO-CI-DI-VA-PLU}

O art. 1º, nos incisos I a V, apresenta os fundamentos da República:

\begin{tabularx}{\linewidth}{llX}
\toprule
 \textbf{Inciso} & \textbf{Fundamento} & \textbf{Ideia central} \\
\midrule
 I & soberania & independência do Estado no plano externo e poder supremo no interno \\
 II & cidadania & participação política e pertencimento à comunidade estatal \\
 III & dignidade da pessoa humana & valor central da ordem constitucional \\
 IV & valores sociais do trabalho e da livre iniciativa & equilíbrio entre trabalho, atividade econômica e justiça social \\
 V & pluralismo político & convivência de diferentes ideias, grupos e projetos políticos \\
\bottomrule
\end{tabularx}


Uma forma de memorização é \textbf{SO-CI-DI-VA-PLU}: soberania, cidadania, dignidade, valores sociais do trabalho e livre iniciativa, pluralismo político.

\subsubsection*{2.1. Soberania}

Soberania não se confunde com autonomia. A República Federativa do Brasil é soberana; Estados, Municípios e Distrito Federal são autônomos nos limites da Constituição. A autonomia dos entes federados não equivale à independência.

\subsubsection*{2.2. Cidadania}

Cidadania é fundamento da República, mas os direitos políticos são apenas uma de suas manifestações. O conceito também envolve participação na vida pública e possibilidade de influenciar as decisões estatais.

\subsubsection*{2.3. Dignidade da pessoa humana}

A dignidade orienta a interpretação dos direitos fundamentais e limita o exercício do poder estatal. Não autoriza, sozinha, qualquer conclusão específica sem análise do caso, mas funciona como vetor interpretativo central.

\subsubsection*{2.4. Trabalho e livre iniciativa}

A Constituição protege tanto os valores sociais do trabalho quanto a livre iniciativa. A ordem constitucional não adota uma visão de liberdade econômica sem responsabilidade social nem elimina a iniciativa privada.

\subsubsection*{2.5. Pluralismo político}

Pluralismo político é mais amplo que a existência de vários partidos. Protege a diversidade de opiniões, ideologias, associações e projetos de sociedade, desde que respeitados os limites constitucionais.

\subsection*{3. Separação dos Poderes}

O art. 2º estabelece que são Poderes da União, independentes e harmônicos entre si, o Legislativo, o Executivo e o Judiciário.

São características complementares:

\begin{itemize}
  \item \textbf{independência:} cada Poder possui funções, organização e garantias próprias;
  \item \textbf{harmonia:} os Poderes devem cooperar e respeitar os limites recíprocos;
  \item \textbf{não exclusividade absoluta:} cada Poder exerce uma função típica, mas também pode exercer funções atípicas previstas na Constituição.
\end{itemize}

O Legislativo tem como função típica legislar e fiscalizar. O Executivo tem como função típica administrar. O Judiciário tem como função típica exercer a jurisdição. Em situações constitucionais, cada um pode exercer funções secundárias: o Executivo edita atos normativos autorizados, o Legislativo administra sua estrutura e o Judiciário organiza seus serviços, por exemplo.

Essa distribuição não significa que um Poder possa substituir livremente o outro. A atuação atípica precisa ter fundamento constitucional ou legal e respeitar controles institucionais.

\subsection*{4. Objetivos fundamentais da República}

O art. 3º apresenta objetivos fundamentais, isto é, metas que orientam a ação do Estado brasileiro:

\begin{enumerate}
  \item construir uma sociedade livre, justa e solidária;
  \item garantir o desenvolvimento nacional;
  \item erradicar a pobreza e a marginalização e reduzir as desigualdades sociais e
\end{enumerate}

   regionais;

\begin{enumerate}
  \item promover o bem de todos, sem preconceitos de origem, raça, sexo, cor, idade
\end{enumerate}

   e quaisquer outras formas de discriminação.

Memorização: \textbf{sociedade, desenvolvimento, pobreza/desigualdades e bem de todos}.

Os objetivos não são apenas declarações políticas sem efeito. Eles orientam a formulação de políticas públicas, a interpretação constitucional e o controle de medidas estatais. Também não devem ser confundidos com os fundamentos do art. 1º.

\subsubsection*{4.1. Erradicar e reduzir}

O texto usa verbos diferentes:

\begin{itemize}
  \item erradicar a pobreza e a marginalização;
  \item reduzir as desigualdades sociais e regionais.
\end{itemize}

Em questão literal, trocar esses verbos pode tornar a alternativa incorreta.

\subsection*{5. Relações internacionais}

O art. 4º lista os princípios que regem o Brasil em suas relações internacionais:

\begin{itemize}
  \item independência nacional;
  \item prevalência dos direitos humanos;
  \item autodeterminação dos povos;
  \item não intervenção;
  \item igualdade entre os Estados;
  \item defesa da paz;
  \item solução pacífica dos conflitos;
  \item repúdio ao terrorismo e ao racismo;
  \item cooperação entre os povos para o progresso da humanidade;
  \item concessão de asilo político.
\end{itemize}

Memorização por blocos:

\begin{enumerate}
  \item \textbf{autonomia:} independência nacional, autodeterminação e não intervenção;
  \item \textbf{igualdade e paz:} igualdade entre Estados, defesa da paz e solução
\end{enumerate}

   pacífica;

\begin{enumerate}
  \item \textbf{direitos humanos:} prevalência dos direitos humanos, repúdio ao
\end{enumerate}

   terrorismo e ao racismo, asilo político;

\begin{enumerate}
  \item \textbf{cooperação:} cooperação entre os povos para o progresso da humanidade.
\end{enumerate}

O parágrafo único do art. 4º determina que a República Federativa do Brasil buscará a integração econômica, política, social e cultural dos povos da América Latina, visando à formação de uma comunidade latino-americana de nações.

\section*{Como a banca cobra}

\begin{itemize}
  \item \textbf{Troca de categoria:} cidadania aparece como objetivo, ou redução das desigualdades aparece como fundamento.
  \item \textbf{Troca de verbo:} a Constituição fala em reduzir desigualdades, não erradicá-las.
  \item \textbf{Federalismo:} Estados e Municípios são autônomos, não soberanos.
  \item \textbf{Poder popular:} o poder emana do povo, não exclusivamente dos eleitos.
  \item \textbf{Lista internacional:} igualdade entre os Estados não significa superioridade de um Estado sobre os demais.
  \item \textbf{Separação:} independência não elimina harmonia; harmonia não elimina controle recíproco.
\end{itemize}

\section*{Exemplos resolvidos}

\subsection*{Exemplo 1}

Um Estado da Federação decide declarar sua independência do Brasil por meio de plebiscito estadual. A medida é compatível com a Constituição?

\textbf{Não.} A República Federativa do Brasil é formada pela união indissolúvel dos Estados, Municípios e Distrito Federal. Os entes federados possuem autonomia, mas não soberania nem direito constitucional de secessão.

\subsection*{Exemplo 2}

Uma alternativa afirma: "A erradicação das desigualdades regionais é um dos fundamentos da República".

\textbf{Incorreta.} A redução das desigualdades sociais e regionais é objetivo fundamental do art. 3º, III. Fundamentos estão no art. 1º e incluem soberania, cidadania, dignidade, valores sociais do trabalho e livre iniciativa e pluralismo político.

\section*{Quadro-resumo}

\begin{tabularx}{\linewidth}{lX}
\toprule
 \textbf{Categoria} & \textbf{Conteúdo principal} \\
\midrule
 Forma de Estado & Federação \\
 Forma de governo & República \\
 Regime político & Democracia \\
 Fundamento central & Dignidade da pessoa humana \\
 Poder & Emana do povo \\
 Poderes & Legislativo, Executivo e Judiciário \\
 Objetivos & sociedade, desenvolvimento, redução das desigualdades e bem de todos \\
 Relações internacionais & direitos humanos, paz, não intervenção e cooperação \\
 Entes federados & autônomos, não soberanos \\
\bottomrule
\end{tabularx}


\section*{Questões de fixação}

\subsection*{Questão 1 — (questão inédita)}

Assinale a alternativa que apresenta apenas fundamentos da República Federativa do Brasil.

A) Soberania, cidadania e dignidade da pessoa humana. \\ B) Defesa da paz, pluralismo político e desenvolvimento nacional. \\ C) Erradicação da pobreza, livre iniciativa e não intervenção. \\ D) Prevalência dos direitos humanos, cidadania e igualdade entre os Estados. \\ E) Construção de sociedade justa, soberania e defesa da paz.

\begin{alerta}
\textbf{Gabarito: A.} Os três elementos estão no art. 1º. As demais alternativas misturam fundamentos com objetivos ou princípios internacionais.
\end{alerta}

\subsection*{Questão 2 — (questão inédita)}

De acordo com a Constituição Federal, o poder:

A) emana exclusivamente dos representantes eleitos. \\ B) emana do povo e pode ser exercido diretamente, nos termos da Constituição. \\ C) pertence aos Poderes da União. \\ D) é exercido apenas por meio de plebiscitos. \\ E) decorre da soberania dos Estados.

\begin{alerta}
\textbf{Gabarito: B.} O parágrafo único do art. 1º prevê exercício por representantes eleitos ou diretamente, conforme a Constituição.
\end{alerta}

\subsection*{Questão 3 — (questão inédita)}

É objetivo fundamental da República Federativa do Brasil:

A) garantir a soberania dos Municípios. \\ B) erradicar as desigualdades regionais. \\ C) reduzir as desigualdades sociais e regionais. \\ D) assegurar a intervenção nos Estados estrangeiros. \\ E) estabelecer uma religião oficial.

\begin{alerta}
\textbf{Gabarito: C.} O art. 3º, III, prevê erradicar a pobreza e a marginalização e reduzir as desigualdades sociais e regionais.
\end{alerta}

\subsection*{Questão 4 — (questão inédita)}

É princípio que rege o Brasil nas relações internacionais:

A) dependência econômica. \\ B) intervenção obrigatória. \\ C) superioridade militar. \\ D) prevalência dos direitos humanos. \\ E) submissão dos Estados menores.

\begin{alerta}
\textbf{Gabarito: D.} A prevalência dos direitos humanos está expressamente no art. 4º, II.
\end{alerta}

\subsection*{Questão 5 — (questão inédita)}

Sobre a separação dos Poderes, a Constituição estabelece que Legislativo, Executivo e Judiciário são:

A) subordinados entre si. \\ B) independentes e harmônicos entre si. \\ C) independentes e sem qualquer forma de controle recíproco. \\ D) hierarquizados conforme a função exercida. \\ E) órgãos exclusivos da União, sem funções atípicas.

\begin{alerta}
\textbf{Gabarito: B.} É a redação literal do art. 2º. Independência e harmonia coexistem, e a Constituição admite o exercício de funções típicas e atípicas.
\end{alerta}

\section*{Revisão final}

\begin{itemize}
  \item[$\square$] Sei diferenciar República, Federação e Estado Democrático de Direito.
  \item[$\square$] Memorizei os cinco fundamentos do art. 1º.
  \item[$\square$] Sei que o poder emana do povo.
  \item[$\square$] Diferencio autonomia de soberania.
  \item[$\square$] Sei os três Poderes e as ideias de independência e harmonia.
  \item[$\square$] Memorizei os quatro objetivos do art. 3º.
  \item[$\square$] Não troco "reduzir desigualdades" por "erradicar desigualdades".
  \item[$\square$] Reconheço os principais princípios internacionais do art. 4º.
\end{itemize}

\end{document}
